(1) 波速 \(c=\sqrt{T/\mu}\)，取 \(y(x,t)=\Re[Y(x)e^{-i\omega t}]\)，则空间满足 \(Y''+k^2Y=0\)，\(k=\omega/c\)，边界条件 \(Y(0)=A\)、\(Y(L)=0\)。这是一类非齐次Dirichlet边值问题，驱动端不是自由端。
(2) 当 \(\sin(kL)\ne0\)，满足边界的唯一同频响应为 \(Y(x)=A\sin[k(L-x)]/\sin(kL)\)，故 \(y(x,t)=A\sin[k(L-x)]\cos(\omega t)/\sin(kL)\)。代入两端分别得到 \(A\cos\omega t\) 和零，代入波动方程也严格成立。
(3) 当 \(kL=n\pi\)，空间齐次问题存在本征模，指定非零边界激励与纯同频有限幅响应通常不兼容，分母趋零。对理想零阻尼且初态有限的共振驱动，受激模振幅在一段时间内近似线性增长，不能把分母零理解为任意时刻具有无限能量。
(4) 弦的局地能流为 \(P_x=-T y_x y_t\)。在无耗非共振的纯实空间同频稳态，驱动端 \(y_x\) 与 \(y_t\) 相差四分之一周期，所以 \(\langle -T y_x(0,t)y_t(0,t)\rangle=0\)，意味着一周期向弦输入的净能量为零。含阻尼时响应有相位差，平均功率可为正；有限时间暂态也会吸收并储存能量。
